% Sanskroot integration blueprint
\documentclass[12pt]{article}
\usepackage[margin=1in]{geometry}
\usepackage{hyperref}
\usepackage{longtable}
\usepackage{amsmath}

% --- Font strategy (XeLaTeX) ---
% This document must be compiled with XeLaTeX to guarantee correct Devanagari rendering.
\usepackage{fontspec}
\defaultfontfeatures{Ligatures=TeX}

% Overleaf includes Google Noto fonts. These two cover Latin + Devanagari reliably.
\setmainfont{Noto Serif}
\setmonofont{Noto Sans Mono}

\newfontfamily\devanagarifont{Noto Sans Devanagari}[Script=Devanagari]
\newcommand{\textdevanagarifont}[1]{{\devanagarifont #1}}

% Automatic font switching for Devanagari code points.
\usepackage{ucharclasses}
\setTransitionsFor{Devanagari}{\devanagarifont}{\normalfont}



\begin{document}

\title{Sanskroot Integration Blueprint for MERKABA Factory}
\author{}
\date{\today}

\maketitle

\tableofcontents

\section{Introduction}

The purpose of this document is to provide a comprehensive blueprint for
integrating \emph{Sanskroot}, a new programming language with source code
written in Devanagari script, into the \textbf{MERKABA Factory} project.  MERKABA
Factory is a deterministic software foundry that converts inputs into
canonical intermediate representations and then synthesises code and other
artefacts.  This blueprint is meant to serve as a ``treasure map'' for
developers and autonomous agents: it outlines the design goals of Sanskroot,
describes the linguistic features of the Devanagari script, explains why
normalisation is necessary, and specifies the implementation tasks required
to ensure that Sanskroot fits within MERKABA’s architecture without
compromising determinism or reproducibility.

The blueprint caters both to minimal viable functionality and to maximal
expressiveness.  It provides a clear specification for a small core language
that is easy to implement, yet lays the groundwork for elaborate features
such as using Devanagari punctuation and special symbols as operators.

\section{MERKABA Factory Overview}

MERKABA Factory exposes a layered architecture with a deterministic
transformation engine at its core.  Runs are described by a run descriptor
and executed by the engine, which applies a pipeline of operators to
generate canonical digests and evidence traces.  On top of the core engine
are domain translators and artefact emitters.  Translators take source
inputs (code or specifications) and produce a \emph{Universal Intermediate
Representation} (UIR), a graph-based structure that captures program
semantics.  Emitters consume this UIR to produce target code, documentation
and deployment bundles.  A vacuum layer manages persistence and reusable
crystals, while surfaces (CLI, server, GUI, WASM) provide user access.

Determinism is enforced by canonicalising inputs, hashing each transformation
step and storing a replayable trace.  For Sanskroot integration, this
architecture implies that the new language must (i) map into the existing
UIR without introducing a new IR, (ii) respect the canonicalisation rules
so that repeated runs over identical inputs produce identical outputs, and
(iii) work with existing emitters or define new emitters in a deterministic
manner.

\section{The Devanagari Script}

\subsection{Script Basics}

Devanagari is an \emph{abugida}, a segmental writing system in which each
consonant letter carries an inherent vowel and additional vowel sounds are
written using dependent marks.  The script is written left-to-right and is
recognisable by a horizontal headline across the top of full letters.  It
descends from the Brāhmī script and is widely used across the Indian
subcontinent, including for Hindi, Marathi and classical Sanskrit.  The
script does not distinguish upper and lower case.  Each
consonant has an inherent vowel \textsl{a}; this vowel is suppressed with
a special sign (halant) and modified with combining marks (matras).  In
tables of consonant–vowel combinations, the bare consonant such as
\textdevanagarifont{क} (\emph{k}) represents the syllable \emph{ka} where the
inherent vowel \textsl{a} is present.

Dependent vowel signs (matras) attach to the base consonant before,
after, above or below.  A vowel sign may have multiple parts that appear on\footnote{Practical notes on Devanagari vowel signs (matras) and storage order: \url{https://r12a.github.io/scripts/deva/hi.html}.}
different sides of the base consonant cluster; these are known as
\emph{matras}.  For example, the short \emph{i}
vowel sign (\devanagarifont{ि}) is stored after the consonant in the
character stream but rendered to the left of the base consonant.  The
underlying code points preserve pronunciation order: the base consonant
character is followed by the vowel sign.  This
character order is critical for deterministic processing and must not be
altered by renderer repositioning.

\subsection{Combining Marks and Normalisation}

Characters in the Devanagari block often involve combining marks such as
vowel signs, nasal marks and the nuqta diacritic.  These marks can occur in
multiple orders that are visually identical; a normalisation procedure is
needed to ensure that equivalent sequences have a unique binary
representation.  The Unicode Standard’s Normalisation Forms define how to\footnote{Unicode Standard Annex \#15 (Normalization Forms), canonical equivalence and NFC: \url{https://unicode.org/reports/tr15/}.}
decompose and reorder combining marks and recombine precomposed characters.
In particular, Normalisation Form C (NFC) performs canonical
decomposition followed by canonical composition.  According to the Unicode
consortium, the normalisation algorithm orders combining marks and uses
decomposition and composition rules so that a binary comparison of
normalised strings determines equivalence.  In other
words, two canonically equivalent Devanagari sequences normalised to NFC
will have exactly the same sequence of code points.

There are script-specific exceptions: certain precomposed Devanagari
characters, such as U+0958 \devanagarifont{क़} (DEVANAGARI LETTER QA), are
excluded from canonical composition.  Implementations
must rely on Unicode normalisation data (e.g. the composition exclusion
tables) to avoid reintroducing such characters during composition.  For
Sanskroot, the canonicalisation phase must normalise input to NFC and
forbid forbidden characters, ensuring that the digest chain remains stable.

\section{Designing the Sanskroot Language}

\subsection{Goals and Philosophy}

Sanskroot is conceived as a playful yet rigorous language written entirely
in Devanagari script.  It demonstrates the flexibility of MERKABA
Factory’s language infrastructure: by adding a new front-end to the
translator layer, one can support a language whose tokens and keywords
belong to a non-Latin script while still benefiting from deterministic
pipelines.  The design intentionally balances minimalism and maximal
capability:

\begin{itemize}
  \item \textbf{Minimal core:} Define a small set of language constructs (function
    definitions, conditionals, variable binding, arithmetic, printing,
    return) that can be mapped to existing UIR nodes.  This allows an
    implementer to create a working prototype quickly.
  \item \textbf{Maximal expressiveness:} Leave room for future extensions such
    as Devanagari punctuation tokens (e.g. danda \devanagarifont{।}, double
    danda \devanagarifont{॥}), special symbols like the om sign
    \devanagarifont{ॐ}, Unicode spaces, or even mixed scripts.  The
    blueprint outlines how to incorporate such features without breaking
    determinism.
\end{itemize}

\subsection{Keyword Vocabulary}

The core language uses a set of Devanagari keywords, chosen both for
readability to speakers of Indic languages and for uniqueness when lexing.
Table~\ref{tab:keywords} lists the keywords together with their Latin
equivalents and meanings.  The keywords are case-insensitive (Devanagari
has no case distinction) and are recognised only after Unicode
normalisation to NFC.

\begin{longtable}{llp{7cm}}
\caption{Core Sanskroot keywords}\label{tab:keywords}\\
\textbf{Devanagari} & \textbf{Latin gloss} & \textbf{Purpose}\\\hline
\devanagarifont{कार्य} & \emph{kārya} & Start of a function definition (analogous to
\texttt{fn}). The keyword is followed by the function name and parameter list.\\
\devanagarifont{यदि} & \emph{yadi} & Conditional branch introduction (\texttt{if}).\\
\devanagarifont{अन्यथा} & \emph{anyathā} & Else branch introduction (\texttt{else}).\\
\devanagarifont{दर्शय} & \emph{darśaya} & Output a value to the console (\texttt{print}).\\
\devanagarifont{निवृत्ति} & \emph{nivṛtti} & Return a value from a function.\\
\devanagarifont{सत्य} & \emph{satya} & Boolean literal true.\\
\devanagarifont{असत्य} & \emph{asatya} & Boolean literal false.\\
\end{longtable}

Additional identifiers (variable names, function names) may use any
sequence of Devanagari letters and combining marks.  The lexer must
normalise to NFC before matching keywords.  Numeric literals follow
international digits by default (\texttt{0–9}); Devanagari digits may be
added later.  String literals are enclosed in double quotation marks
(\texttt{"} and \devanagarifont{"}) and may contain arbitrary Unicode
characters; escapes follow Rust/Python conventions.

\subsection{Lexical Rules}

\paragraph{Normalisation.} Before lexing, the entire input must be
normalised to Unicode NFC.  Without this, the same visual source could have
different byte sequences (e.g. combining marks in varying order).  NFC
normalises combining marks and recomposes sequences; it ensures that
canonically equivalent texts reduce to a unique form.

\paragraph{Identifiers.} An identifier begins with a Devanagari letter and may
be followed by Devanagari letters, combining marks (matras, nasalisation
marks), the nuqta (\devanagarifont{़}), and underscores.  Formally, let
\texttt{XID\_Start\_Deva} denote the set of Devanagari code points of
general category \texttt{Lu}, \texttt{Ll} and \texttt{Lo}.  Then an
identifier token matches

\[\text{identifier} := \mathrm{XID\_Start\_Deva}\ (\mathrm{XID\_Continue\_Deva})^*\]

where \texttt{XID\_Continue\_Deva} includes letters, combining marks and
digits.  The lexer checks if the resulting NFC-normalised token matches a
keyword; otherwise it is treated as an identifier.

\paragraph{Whitespace and Comments.} Standard ASCII whitespace (space, tab,
newline) separates tokens.  As an extension, the Devanagari danda
(\devanagarifont{।}) and double danda (\devanagarifont{॥}) may serve as
statement separators.  Single-line comments start with a hash (\texttt{\#})
and extend to the end of the line.  Multiline comments use
\texttt{/*} \dots \texttt{*/}.

\paragraph{Operators.} For the minimal version, operators mirror those of a
typical imperative language: arithmetic (+, –, *, /), comparison (==, !=,
<, >, \texttt{≤}, \texttt{≥}), logical (\texttt{\&\&}, \texttt{||}), and
assignment (=).  Operators are ASCII for ease of typing.  The extended
version may define Devanagari punctuation as operators (see
Section~\ref{sec:maximum}).

\subsection{Grammar (EBNF)}

The core grammar of Sanskroot is intentionally simple and is presented in
Extended Backus–Naur Form (EBNF).  Terminals are highlighted in
\devanagarifont{Devanagari} where appropriate.

\begin{verbatim}
program       ::= { function }

function      ::= 'कार्य' identifier '(' [ paramlist ] ')' '{' { statement } '}'
paramlist     ::= identifier { ',' identifier }
statement     ::= block | conditional | return | print | assignment | expression ';'
block         ::= '{' { statement } '}'
conditional   ::= 'यदि' '(' expression ')' statement [ 'अन्यथा' statement ]
return        ::= 'निवृत्ति' expression ';'
print         ::= 'दर्शय' expression ';'
assignment    ::= identifier '=' expression ';'
expression    ::= logical { ('+' | '-' ) logical }
logical       ::= comparison { ('*' | '/') comparison }
comparison    ::= unary { ( '==' | '!=' | '<' | '>' | '<=' | '>=' ) unary }
unary         ::= primary | ('-' | '!') unary
primary       ::= boolean | number | string | identifier | '(' expression ')'
boolean       ::= 'सत्य' | 'असत्य'
number        ::= [0-9]+  ; future: Devanagari digits
string        ::= '"' { any-character-but-" } '"'
\end{verbatim}

This grammar yields an abstract syntax tree (AST) with nodes: Program,
Function, Block, IfElse, Return, Print, Assignment, BinaryOp, UnaryOp,
Literal, Identifier.  When mapping to the UIR, each AST node becomes one or
more UIR nodes (described below).

\subsection{UIR Mapping}

MERKABA’s domain translator defines a graph-based UIR with node types such
as Module, Function, Statement, Expression, Literal, Symbol and edges to
capture containment, call relationships and data flow.  Sanskroot’s AST can
be mapped to UIR as follows:

\begin{itemize}
  \item \textbf{Program}: Create a \texttt{Module} node labelled with the
    file name or module name.  Each function becomes a child via a
    \texttt{Contains} edge.
  \item \textbf{Function}: Create a \texttt{Function} node with name and
    parameter list.  Parameter identifiers become \texttt{Symbol} nodes with
    a \texttt{Defines} edge from the function.  The function’s body is
    represented as a \texttt{Block} node.
  \item \textbf{Conditional}: Create an \texttt{If} node.  Connect the
    condition expression to the \texttt{If} node via a \texttt{Condition}
    edge.  The consequent and optional alternative blocks are connected via
    \texttt{Then} and \texttt{Else} edges.
  \item \textbf{Return}: Create a \texttt{Return} node and link the return
    expression via an \texttt{Expr} edge.  Connect the return node to the
    surrounding function via a \texttt{Contains} edge.
  \item \textbf{Print}: Create a \texttt{Call} node targeting a built-in
    print function.  The argument expression is attached via an
    \texttt{Argument} edge.
  \item \textbf{Assignment}: Create a \texttt{Defines} edge from the current
    scope to a \texttt{Symbol} node representing the variable.  Map the
    right-hand side expression to UIR and connect via an \texttt{Init}
    edge.
  \item \textbf{Expression, BinaryOp, UnaryOp, Literal}: Use existing UIR
    expression nodes with appropriate operator labels.  Binary operations
    become \texttt{BinaryExpr} nodes with \texttt{Left} and \texttt{Right}
    edges.  Literals become \texttt{Literal} nodes.
\end{itemize}

This mapping strategy leverages the existing UIR without introducing a new
representation.  Since MERKABA already serialises UIR via a DTO, the
Sanskroot translator must produce UIR in the canonical DTO form (e.g.
\texttt{UirGraphDto}).

\subsection{Implementation Guidance}

Implementing Sanskroot requires writing a new frontend in the MERKABA
repository, for example in \texttt{kernel/crates/domain-translator/src/frontends/sanskroot.rs}.
Key steps include:

\begin{enumerate}
  \item \textbf{Normalise input.} Use a Unicode normalisation library
    (e.g. Rust’s \texttt{unicode\_normalization}) to convert the source into
    NFC before any lexing.  This ensures that visually identical sources
    yield identical byte sequences and thus identical digests.

  \item \textbf{Lexing.} Implement a lexer that recognises Devanagari
    keywords, identifiers, numeric literals and string literals.  Keywords
    should be compared on NFC-normalised strings.  Record token positions
    for error reporting.

  \item \textbf{Parsing.} Implement a recursive descent or Pratt parser
    matching the EBNF grammar.  Build an AST representing the program.

  \item \textbf{Lowering to UIR.} Traverse the AST and emit UIR nodes and
    edges according to the mapping described above.  Ensure that UIR node
    identifiers are deterministic (e.g. derived from input order) so that
    repeated runs generate the same UIR structure.

  \item \textbf{Integration.} Expose Sanskroot as a supported language in
    the translator’s configuration (e.g. via a language registry).  Add
    tests to the repository that compile sample Sanskroot programs,
    generate UIR JSON and verify that repeated runs produce identical
    digests.  Provide CLI flags so that users can specify \texttt{--source-lang
    sanskroot}.

  \item \textbf{Documentation.} Update user documentation and the README to
    describe the new language, including the keywords and grammar.  Include
    examples of round-trip translation to existing targets (e.g. Python or
    Rust code generation).
\end{enumerate}

\subsection{Definition of Done}

To ensure that the Sanskroot integration is successful, the following
criteria must be met:

\begin{enumerate}
  \item \textbf{Deterministic replay.} Two runs over the same Sanskroot
    program and run descriptor must yield the same output digests and UIR
    JSON.  This implies that the parser and UIR emission are deterministic
    and that NFC normalisation is applied consistently.

  \item \textbf{Canonical UIR.} The output UIR must use the canonical
    transport format adopted by MERKABA (e.g. \texttt{UirGraphDto}).  No
    additional intermediate representations may be introduced.

  \item \textbf{Golden path.} A CLI command like:
    \begin{quote}
      \texttt{merkaba run --source-lang sanskroot --input examples/hello.skrt \\
      --emit-bundle out/hello\_bundle}
    \end{quote}
    must generate a bundle containing the normalised source, UIR JSON,
    evidence trace, manifest, and generated code (e.g. Python).  A test
    should verify that the bundle is identical across repeated runs.

  \item \textbf{Tests and documentation.} Unit tests must cover the
    lexer, parser and UIR mapping for Sanskroot.  The documentation must
    accurately reflect the implemented behaviour.

  \item \textbf{Compatibility.} The integration must not break existing
    languages or the core MERKABA workflow.  Code review should confirm
    that common utilities (e.g. normalisation, digest policy) are used
    consistently.
\end{enumerate}

\section{Maximum Features and Extensions}
\label{sec:maximum}

This section sketches optional extensions that push Sanskroot into
\emph{maximal} territory.  Implementers should treat these as future work
and only pursue them after the minimal core is stable.

\subsection{Devanagari Punctuation as Operators}

The Devanagari script includes punctuation such as the danda (\devanagarifont{।})
and double danda (\devanagarifont{॥}).  These symbols can be repurposed as
statement separators or block delimiters.  For example, instead of using
curly braces, one might use \devanagarifont{॥} to close a block.  If this
extension is adopted, the parser must normalise these punctuation marks and
define clear operator precedence.  Any such feature should come with a
normalisation rule to avoid confusable characters (e.g. full stop vs. danda).

\subsection{Om Sign for Entry Point}

The Om sign (\devanagarifont{ॐ}) is a visually distinctive character with deep
symbolic meaning.  A maximal variant of Sanskroot could designate a
function whose name is \devanagarifont{ॐ} as the program’s entry point, similar
to \texttt{main} in C.  The translator would interpret a top-level
\texttt{कार्यॐ()} definition as the entry point.  Because this symbol is
excluded from canonical composition, NFC normalisation preserves its code
point.

\subsection{Mixed Scripts and Unicode Fun}

One might allow ASCII identifiers or operators inside a Devanagari program,
or even embed emojis.  To maintain determinism, any such extension must
define a precise normalisation pipeline and a comprehensive keyword table.
Unicode confusable characters should be mapped to a canonical set to avoid
hash differences.  Implementers are advised to restrict mixed scripts to
well-defined contexts (e.g. string literals) until strong normalisation
support is available.

\subsection{Custom Digits and Numerals}

Devanagari digits (\devanagarifont{०}–\devanagarifont{९}) could be supported for
numeric literals.  To do this, the lexer must translate these symbols to
their corresponding numeric values and ensure that normalisation does not
alter them.  The Unicode standard includes composition exclusions for
certain Devanagari characters, but digits are
composed and decompose trivially.

\section{Integration Checklist}

For ease of reference, this section summarises the critical tasks that must
be completed to integrate Sanskroot into MERKABA Factory:

\begin{enumerate}
  \item \textbf{Verify Devanagari support.} Ensure that build tools and
    continuous integration servers support Unicode source files.  Use
    version control settings (e.g. \texttt{.gitattributes}) to treat
    Sanskroot files as UTF-8 text.
  \item \textbf{Normalisation.} Incorporate NFC normalisation into the
    canonicalisation phase of the MERKABA pipeline.  Write unit tests
    demonstrating that multiple canonical Devanagari encodings map to the
    same digest.
  \item \textbf{Frontend implementation.} Implement the lexer, parser,
    AST and UIR mapping as described.  Use deterministic node ordering
    and avoid non-deterministic data structures (e.g. hash maps without
    deterministic iteration).
  \item \textbf{Bundle generation.} Extend the CLI and server surfaces
    to accept Sanskroot sources and to emit bundles containing the
    translated code and metadata.  Update the GUI (Tauri) to support
    Devanagari file names and syntax highlighting.
  \item \textbf{Testing and review.} Add golden path tests and replay
    verification.  Conduct code review to ensure that no global state or
    randomisation has crept into the implementation.
\end{enumerate}

\section{Conclusion}

This blueprint lays the foundation for integrating Sanskroot, a
Devanagari-based programming language, into MERKABA Factory.  By adhering
to Unicode normalisation, mapping to the existing UIR and maintaining
determinism throughout the pipeline, developers can build a rich and
playful language without jeopardising the core guarantees of MERKABA.
Once the minimal language is fully implemented and verified, the maximal
extensions sketched here—including Devanagari punctuation operators and
special entry symbols—offer fertile ground for experimentation and
trolling, turning the system into a versatile and fantastical programming
kitchen.

\end{document}
